\section{Conclusion}\label{sec:concl}

We presented \sname, a LiDAR localization stack whose contribution is a specific integration of established components rather than a new algorithm: three known estimator families, likelihood-field MCL, NDT-MCL \citep{saarinen2013ndtmcl}, and an error-state Kalman filter fusing LiDAR, IMU, and RTK-GNSS, sit behind a single ROS~2 topic and TF contract, so that a 2D LiDAR, a 3D LiDAR, or a sensor-fused configuration is a deployment choice rather than a switch to a stack with a different interface. Each estimator is a pure C++ core with no dependence on the middleware runtime, the two MCL backends share one predict--weight--resample core that differs only in its measurement model, as in existing MCL libraries \citep{puga2024beluga}, and \code{laser2d} relocalizes deterministically by branch-and-bound \citep{hess2016realtime} over a map-sized window of non-occupied cells and verifies the top candidates over several scans before committing, a simple instance of established multi-hypothesis verification that can report a place as ambiguous. Seeded synthetic experiments show that the assembled pipelines behave as intended in closed loop, and they also expose two open problems that we report rather than hide: on large maps with repeated structure a single scan admits wrong poses, which multi-scan verification makes rare (1 of 240 synthetic queries) but does not rule out, and the ESKF covariance is overconfident, mainly in attitude. An evaluation on real and public data against external baselines remains future work. \sname{} is released under the Apache-2.0 license at \url{https://github.com/kjungmo/prism_loc}.
